\documentclass{tsp-short}
\usepackage{graphicx}
\usepackage{amsthm}
\newtheorem{assumption}{Assumption}
%%
%--%Place your definitions here
\newtheorem{theorem}{Theorem}[section]
\newtheorem{lemma}{Lemma}[section]
\newtheorem{corollary}{Corollary}[section]
\newtheorem{proposition}{Proposition}[section]
\theoremstyle{remark}
\newtheorem{remark}{Remark}[section]
\theoremstyle{definition}
\newtheorem{definition}{Definition}[section]
\newtheorem{example}{Example}[section]
\newcommand{\Ent}{\mathrm{Ent}}

\title[Interacting SDEs]{Stochastic Differential Equations with Interaction for Non-Constant Diffusion}

\author{Kyrylo Kuchynskyi}
\address{Ukraine, Kyiv-4, 3, Tereschenkivska st.}
\email{kkuchynskyi@imath.kiev.ua}
\thanks{This work was supported by a grant from the Simons Foundation (1290607, K.O.K.).}

\subjclass[2020]{60H10, 60B05, 60G57}
\keywords{common noise, interacting SDEs, stochastic flows, logarithmic Sobolev inequality, Jacobian matrix flow}

\begin{document}

\begin{abstract}
We study interacting SDEs driven by common noise where the diffusion coefficient is state-dependent.
We derive the SDE for the Jacobian matrix of the associated stochastic flow and obtain a pathwise exponential bound on directional derivatives, which yields a pathwise logarithmic Sobolev inequality for the transported measures with a random constant controlled by the flow Lipschitz constant.
\end{abstract}

\maketitle


% =========================================================
\section{Introduction}\label{sec:intro}
% =========================================================

We consider families $(\mu_t)_{t\ge 0}$ of random probability measures on $\mathbb{R}^d$ generated by a continuum of particles labeled by their initial position $u\in\mathbb{R}^d$ and driven by the same $d$-dimensional Brownian motion $W_t$:
\begin{equation}
\begin{cases}
dX_t(u) = a\big(X_t(u),\mu_t\big)\,dt + \sigma\big(X_t(u)\big)\,dW_t,\\[2mm]
X_0(u)=u,\\[1mm]
\mu_t = \mu_0\circ X_t^{-1}.
\end{cases}
\label{eq:main_eqn}
\end{equation}
Here $\mu_0$ is a deterministic initial probability measure, the drift $a:\mathbb{R}^d\times\mathcal{P}_2(\mathbb{R}^d)\to\mathbb{R}^d$ may depend on the current measure, and the diffusion coefficient $\sigma:\mathbb{R}^d\to\mathbb{R}^{d\times d}$ is state-dependent and non-constant.



The special feature of \eqref{eq:main_eqn} is the common noise: for each realization $\omega$, the map
\begin{equation*}
T_t(\omega,\cdot):u\mapsto X_t(u,\omega)
\end{equation*}
acts on the initial measure by pushforward,
\begin{equation*}
\mu_t(\omega)=\mu_0\circ T_t(\omega)^{-1},
\end{equation*}
so the evolution is best viewed as a random transport of measures by a stochastic flow.  Such equations are called Stochastic Differential Equations with Interactions and were proposed by Andrey A. Dorogovtsev in \cite{dorogovtsev2003}. 

\begin{assumption}[A2]\label{ass:general_lipschitz}
The coefficients satisfy the following conditions:
\begin{enumerate}
\item[(i)] The map $x\mapsto a(x,\mu)$ is $C^1$ for each $\mu\in\mathcal{P}_2(\mathbb{R}^d)$ and
\begin{equation*}
\sup_{x,\mu}\|\partial_x a(x,\mu)\|_{\mathrm{op}}\le L.
\end{equation*}
\item[(ii)] The diffusion $\sigma:\mathbb{R}^d\to\mathbb{R}^{d\times d}$ is $C^1$ and there exists $L_\sigma<\infty$ such that
\begin{equation*}
\sup_x \|\partial \sigma(x)\|_{\mathrm{op}}\le L_\sigma.
\end{equation*}
\item[(iii)] Under (i)--(ii), \eqref{eq:main_eqn} admits a unique global strong solution for each initial condition,
and it generates a $C^1$ stochastic flow of maps: there exists a modification such that for a.s.\ $\omega$ and every $t\ge 0$ the map
$u\mapsto X_t(u,\omega)$ is of class $C^1(\mathbb{R}^d)$ (in particular, it is locally Lipschitz; moreover, it is Lipschitz on every compact set $K\subset\mathbb{R}^d$ with Lipschitz constant $\sup_{u\in K}\|J_t(u,\omega)\|_{\mathrm{op}}$).
We emphasize that such assumptions do not, in general, imply a finite global Lipschitz constant $\sup_{u\in\mathbb{R}^d}\|J_t(u,\omega)\|_{\mathrm{op}}$ for fixed $(t,\omega)$; see, e.g., Mohammed--Scheutzow~\cite{mohammed_scheutzow1998}.
\end{enumerate}
\end{assumption}

\begin{definition}[Logarithmic Sobolev inequality]
Denoting by $\mathcal{C}^1_c(\mathbb{R}^d)$ the set of compactly supported smooth functions from $\mathbb{R}^d$ to $\mathbb{R}$, for $g\ge 0$ with $\int g\,d\mu=1$ we define the entropy
\begin{equation*}
\Ent_\mu(g):=\int g\log g\,d\mu.
\end{equation*}
A probability measure $\mu$ on $\mathbb{R}^d$ is said to satisfy a logarithmic Sobolev inequality with constant $C>0$ if for every $h\in\mathcal{C}^1_c(\mathbb{R}^d)$ with $\int h^2\,d\mu=1$ one has
\begin{equation}
\Ent_\mu(h^2)\le C\int_{\mathbb{R}^d}|\nabla h|^2\,d\mu.
\label{eq:lsi_def}
\end{equation}
\end{definition}


In the constant diffusion case $\sigma(x)\equiv \sigma_0$, the flow $T_t(\omega)$ is (pathwise) Lipschitz with a deterministic bound with $\mathrm{Lip}(T_t)\le e^{Lt}$ if $a(\cdot,\mu)$ is $L$-Lipschitz in space, which immediately implies propagation of LSI with a deterministic constant $C_0 e^{2Lt}$.

\subsection*{Notation}
We write $|\cdot|$ for the Euclidean norm on $\mathbb{R}^d$.
For a matrix $M\in\mathbb{R}^{d\times d}$ we set
\begin{equation*}
\|M\|_{\mathrm{op}}:=\sup_{|v|=1}|Mv|,
\qquad
\|M\|_{\mathrm{HS}}:=\Big(\sum_{i,j=1}^d M_{ij}^2\Big)^{1/2}.
\end{equation*}
For $\sigma=(\sigma^{(1)},\ldots,\sigma^{(d)})$ with columns $\sigma^{(k)}:\mathbb{R}^d\to\mathbb{R}^d$, we define
\begin{equation*}
\|\partial\sigma(x)\|_{\mathrm{op}}:=\max_{1\le k\le d}\|\partial \sigma^{(k)}(x)\|_{\mathrm{op}}.
\end{equation*}

Given a map $f:\mathbb{R}^d\to\mathbb{R}^m$ of class $C^1$, we write $\partial f(x)$ for its Jacobian matrix, i.e.
$\big(\partial f(x)\big)_{ij}:=\partial_{x_j} f_i(x)$.












% =========================================================





% =========================================================





\section{Control of $\mathrm{Lip}(T_t)$ and LSI pushforwards}\label{sec:jacobian}
% =========================================================

\subsection{Preliminaries: LSI and Lipschitz pushforwards}

\begin{lemma}[LSI stability under Lipschitz pushforwards]\label{lem:lsi-pushforward}
Let $\mu$ satisfy $\mathrm{LSI}(C)$ and let $T:\mathbb{R}^d\to\mathbb{R}^d$ be $K$-Lipschitz. Then $T_\#\mu:=\mu\circ T^{-1}$ satisfies $\mathrm{LSI}(CK^2)$.
\end{lemma}

\begin{proof}
Let $g\in C_c^1(\mathbb{R}^d)$ and set $h=g\circ T$. Then by the change of variables identity,
$\mathrm{Ent}_{T_\#\mu}(g^2)=\mathrm{Ent}_\mu(h^2)$.
Applying \eqref{eq:lsi_def} to $h$ gives
\begin{equation*}
\mathrm{Ent}_{T_\#\mu}(g^2) \le C\int |\nabla(g\circ T)|^2\,d\mu.
\end{equation*}
By the chain rule, $|\nabla(g\circ T)(x)|\le K|\nabla g(T(x))|$ a.e., hence
\begin{equation*}
\int |\nabla(g\circ T)|^2\,d\mu \le K^2\int |\nabla g|^2\,d(T_\#\mu),
\end{equation*}
which yields $\mathrm{LSI}(CK^2)$ for $T_\#\mu$.
\end{proof}

\begin{corollary}[Pathwise LSI bound for \eqref{eq:main_eqn}]\label{cor:pathwise-lsi}
Assume $\mu_0$ satisfies $\mathrm{LSI}(C_0)$. For each $t\ge 0$ and a.s. $\omega$,
\begin{equation*}
\mu_t(\omega)=\mu_0\circ T_t(\omega)^{-1}\ \ \Longrightarrow\ \ 
\mu_t(\omega)\text{ satisfies }\mathrm{LSI}\Big(C_0\,\mathrm{Lip}(T_t(\omega))^2\Big).
\end{equation*}
\end{corollary}

Thus, it remains to control the flow Lipschitz constant $\mathrm{Lip}(T_t(\omega))$.
In general, global (in space) almost sure bounds on spatial derivatives of stochastic flows may fail even under globally Lipschitz assumptions; see \cite{mohammed_scheutzow1998} for counterexamples.

% ---------------------------------------------------------
\subsection{From the Jacobian matrix to $\mathrm{Lip}(T_t)$}\label{subsec:lip-control}
% ---------------------------------------------------------

Fix $(t,\omega)$ and assume that the flow map $u\mapsto T_t(\omega,u):=X_t(u,\omega)$ is of class $C^1$.
We denote its Jacobian matrix by $J_t(u,\omega):=\partial_u T_t(\omega,u)\in\mathbb{R}^{d\times d}$.

\begin{lemma}[Flow Lipschitz constant via Jacobian matrix]\label{lem:lip-via-J}
For each fixed $(t,\omega)$,
\begin{equation*}
\mathrm{Lip}(T_t(\omega)) = \sup_{u\in\mathbb{R}^d}\|J_t(u,\omega)\|_{\mathrm{op}}\in[0,\infty].
\end{equation*}
\end{lemma}

\begin{proof}
Since $u\mapsto T_t(\omega,u)$ is $C^1$, the mean value theorem implies
\begin{equation*}
|T_t(u)-T_t(v)| \le \sup_{\theta\in[0,1]}\|J_t(v+\theta(u-v))\|_{\mathrm{op}}\,|u-v|.
\end{equation*}
Taking the supremum over $u\neq v$ yields $\mathrm{Lip}(T_t)\le \sup_u\|J_t(u)\|_{\mathrm{op}}$.
Conversely, for any $u$ and unit vector $\xi$, letting $v=u+\varepsilon\xi$ and sending $\varepsilon\downarrow 0$ gives
\begin{equation*}
\|J_t(u)\|_{\mathrm{op}} = \sup_{|\xi|=1}\lim_{\varepsilon\downarrow 0}\frac{|T_t(u+\varepsilon\xi)-T_t(u)|}{\varepsilon}
\le \mathrm{Lip}(T_t),
\end{equation*}
hence equality.
\end{proof}










% =========================================================





\section{Jacobian matrix theorem}\label{sec:jacobian-theorem}
% ---------------------------------------------------------

For a fixed $\omega$ and $t$, define the flow map $T_t(\omega,u):=X_t(u,\omega)$ and its Jacobian matrix
\begin{equation*}
J_t(u):=\partial_u T_t(u)=\partial_u X_t(u)\in\mathbb{R}^{d\times d}.
\end{equation*}
Under Assumption~\ref{ass:general_lipschitz}, standard results on stochastic flows (see, e.g., \cite[Theorem~1.4.1]{Kunita1986}) yield a modification such that for a.s.\ $\omega$ and each $t\ge0$ the map $u\mapsto T_t(\omega,u)$ is $\mathcal{C}^1$.
We now state the SDE satisfied by $J_t(u)$.

\begin{theorem}[Jacobian matrix SDE for \eqref{eq:main_eqn}]\label{thm:jacobianSDE}
Under Assumption~\ref{ass:general_lipschitz}, for each $u\in\mathbb{R}^d$, the Jacobian matrix process $J_t(u)$ exists and satisfies the linear matrix SDE (cf.\ \cite{Kunita1986,dorogovtsev2023intermittency})
\begin{equation}
\boxed{
dJ_t(u)= A_t(u)\,J_t(u)\,dt + \sum_{k=1}^d B_t^{(k)}(u)\,J_t(u)\,dW_t^{k},\qquad J_0(u)=I,
}
\label{eq:JacSDE}
\end{equation}
where
\begin{equation*}
A_t(u):=\partial_x a\big(X_t(u),\mu_t\big),\qquad
B_t^{(k)}(u):=\partial_x \sigma^{(k)}\big(X_t(u)\big),
\end{equation*}
and $\sigma^{(k)}$ denotes the $k$-th column of $\sigma$.
Here $A_t(u),B_t^{(k)}(u)\in\mathbb{R}^{d\times d}$, the driving Brownian motion is $W_t=(W_t^1,\ldots,W_t^d)$ with $W_t^k\in\mathbb{R}$, and the sum runs over the $d$ components of the noise.
Moreover, the coefficients satisfy the uniform bounds
\begin{equation*}
\|A_t(u)\|_{\mathrm{op}}\le L,\qquad
\|B_t^{(k)}(u)\|_{\mathrm{op}}\le L_\sigma
\quad \text{for all }t,u,\omega.
\end{equation*}
\end{theorem}

\begin{remark}[One-dimensional picture]
Consider the scalar case $d=1$, where
\begin{equation*}
dX_t(u)=a\big(X_t(u),\mu_t\big)\,dt+\sigma\big(X_t(u)\big)\,dW_t,\qquad X_0(u)=u.
\end{equation*}
Then $J_t(u)=\partial_u X_t(u)$ is scalar and satisfies the linear SDE
\begin{equation*}
dJ_t(u)=\partial_x a\big(X_t(u),\mu_t\big)\,J_t(u)\,dt+\sigma'\big(X_t(u)\big)\,J_t(u)\,dW_t,\qquad J_0(u)=1.
\end{equation*}
This equation can be solved explicitly:
\begin{equation*}
J_t(u)=\exp\Big(\int_0^t\Big(\partial_x a\big(X_s(u),\mu_s\big)-\tfrac12|\sigma'(X_s(u))|^2\Big)\,ds+\int_0^t \sigma'(X_s(u))\,dW_s\Big),
\end{equation*}
so $|J_t(u)|$ is a random exponential factor.
\end{remark}

\subsection{Proof of Theorem~\ref{thm:jacobianSDE}}\label{subsec:jacobian-proof}

\begin{proof}
We give a derivation which makes clear why no derivative of $\mu_t$ appears.

Fix $u\in\mathbb{R}^d$ and a direction $e\in\mathbb{R}^d$ with $|e|=1$. For $h\neq 0$, define the difference quotient
\begin{equation*}
\Delta_t^{(h)} := \frac{X_t(u+he)-X_t(u)}{h}.
\end{equation*}
Subtract the two SDEs in \eqref{eq:main_eqn} (note that both trajectories use the same random measure $\mu_t$ at time $t$ because they belong to the same flow), obtaining
\begin{align*}
d\big(X_t(u+he)-X_t(u)\big)
&= \Big(a(X_t(u+he),\mu_t)-a(X_t(u),\mu_t)\Big)\,dt\\
&\quad + \Big(\sigma(X_t(u+he))-\sigma(X_t(u))\Big)\,dW_t.
\end{align*}

Divide by $h$:
\begin{equation}
d\Delta_t^{(h)}
= \frac{a(X_t(u+he),\mu_t)-a(X_t(u),\mu_t)}{h}\,dt
+ \frac{\sigma(X_t(u+he))-\sigma(X_t(u))}{h}\,dW_t.
\label{eq:diffquot}
\end{equation}

By the mean value theorem and the regularity in Assumption~\ref{ass:general_lipschitz},
\begin{equation*}
\frac{a(X_t(u+he),\mu_t)-a(X_t(u),\mu_t)}{h}
= \int_0^1 \partial_x a\big(X_t(u)+\theta(X_t(u+he)-X_t(u)),\mu_t\big)\,d\theta\ \Delta_t^{(h)},
\end{equation*}
and similarly for $\sigma$, with $\partial\sigma^{(k)}$.
Denote the averaged matrices by $A_t^{(h)}$ and $B_t^{(h,k)}$. Then \eqref{eq:diffquot} is a linear SDE:
\begin{equation*}
d\Delta_t^{(h)} = A_t^{(h)}\,\Delta_t^{(h)}\,dt + \sum_{k=1}^d B_t^{(h,k)}\,\Delta_t^{(h)}\,dW_t^k,
\qquad \Delta_0^{(h)} = e.
\end{equation*}
Uniform boundedness of $A_t^{(h)}$ and $B_t^{(h,k)}$ follows from Assumption~\ref{ass:general_lipschitz}.
Standard stability results for linear SDEs with uniformly bounded coefficients imply that, as $h\to 0$,
$\Delta_t^{(h)}$ converges in probability (and, along subsequences, almost surely) to the directional derivative $J_t(u)e$.
Passing to the limit in the coefficients yields
\begin{equation*}
d(J_t(u)e)= A_t(u)\,(J_t(u)e)\,dt + \sum_{k=1}^d B_t^{(k)}(u)\,(J_t(u)e)\,dW_t^k,
\qquad J_0(u)e=e.
\end{equation*}
Since this holds for every basis vector $e$, it yields the matrix SDE \eqref{eq:JacSDE}.
\end{proof}

% \begin{proof}
% From \eqref{eq:directional-pathwise},
% $|J_t(u)\xi|^{2q}\le \exp(q(2L+dL_\sigma^2)t)\exp(qN_t)$.
% Since $\langle N\rangle_t\le 4dL_\sigma^2 t$, the exponential martingale inequality gives
% \begin{equation*}
% \mathbb{E}\,e^{qN_t}
% =\mathbb{E}\Big[e^{qN_t-\frac12 q^2\langle N\rangle_t}\,e^{\frac12 q^2\langle N\rangle_t}\Big]
% \le \exp\Big(\tfrac12 q^2\cdot 4dL_\sigma^2 t\Big),
% \end{equation*}
% which yields \eqref{eq:moment-direction}.
% For \eqref{eq:moment-op}, note $\|J\|_{\mathrm{op}}\le \|J\|_{\mathrm{HS}}$ and
% $\|J\|_{\mathrm{HS}}^2=\sum_{i=1}^d |Je_i|^2$.
% For $q\ge1$, $(\sum_{i=1}^d a_i)^q\le d^{q-1}\sum_{i=1}^d a_i^q$ for $a_i\ge0$; apply this with $a_i=|J_te_i|^2$
% and then use \eqref{eq:moment-direction} for each $e_i$.
% \end{proof}


\section{Pathwise control of $\|J_t\|$ and a ``Jacobian matrix/Lipschitz constant''}\label{subsec:jacobian-constant}
% ---------------------------------------------------------

Throughout, fix $u\in\mathbb{R}^d$ and a unit vector $\xi\in\mathbb{R}^d$.
Set $Y_t:=J_t(u)\xi$ and $V_t:=|Y_t|^2$.

\begin{proposition}[Directional pathwise exponential bound]\label{prop:directional-exp-bound}
Let $Y_t$ solve
\begin{equation*}
dY_t = A_t(u)Y_t\,dt + \sum_{k=1}^d B_t^{(k)}(u)Y_t\,dW_t^k,
\qquad Y_0=\xi,
\end{equation*}
with $\|A_t(u)\|_{\mathrm{op}}\le L$ and $\|B_t^{(k)}(u)\|_{\mathrm{op}}\le L_\sigma$.
Then for every $t\ge0$,
\begin{equation}\label{eq:directional-pathwise}
|J_t(u)\xi|^2 \le \exp\Big((2L+dL_\sigma^2)t + N_t(u,\xi)\Big)\quad\text{a.s.},
\end{equation}
where $N_t(u,\xi)$ is a continuous local martingale of the form
\begin{equation*}
N_t(u,\xi)=\sum_{k=1}^d\int_0^t \theta_s^{(k)}(u,\xi)\,dW_s^k,
\qquad |\theta_s^{(k)}(u,\xi)|\le 2L_\sigma,
\end{equation*}
hence $\langle N(u,\xi)\rangle_t\le 4dL_\sigma^2 t$.
\end{proposition}

\begin{proof}
By It\^o's formula,
\begin{equation*}
dV_t
=\Big(2\langle Y_t,A_tY_t\rangle+\sum_{k=1}^d|B_t^{(k)}Y_t|^2\Big)dt
+2\sum_{k=1}^d\langle Y_t,B_t^{(k)}Y_t\rangle\,dW_t^k.
\end{equation*}
Using $\langle Y,AY\rangle\le \|A\|_{\mathrm{op}}|Y|^2$ and $|BY|^2\le \|B\|_{\mathrm{op}}^2|Y|^2$ gives
\begin{equation*}
\alpha_t:=2\langle Y_t,A_tY_t\rangle+\sum_{k=1}^d|B_t^{(k)}Y_t|^2 \le (2L+dL_\sigma^2)V_t.
\end{equation*}
For $\varepsilon>0$ apply It\^o to $\log(V_t+\varepsilon)$:
\begin{equation*}
d\log(V_t+\varepsilon)=\frac{dV_t}{V_t+\varepsilon}
-\frac{1}{2}\frac{d\langle V\rangle_t}{(V_t+\varepsilon)^2}
\le (2L+dL_\sigma^2)\,dt+\sum_{k=1}^d\frac{2\langle Y_t,B_t^{(k)}Y_t\rangle}{V_t+\varepsilon}\,dW_t^k,
\end{equation*}
since the quadratic variation term is nonpositive.
Moreover,
$\big|\frac{2\langle Y,B^{(k)}Y\rangle}{V+\varepsilon}\big|\le 2\|B^{(k)}\|_{\mathrm{op}}\le 2L_\sigma$,
so the stochastic integral is well-defined.
Integrating and letting $\varepsilon\downarrow0$ yields \eqref{eq:directional-pathwise} with
$\theta_s^{(k)}(u,\xi):= 2\langle Y_s,B_s^{(k)}Y_s\rangle/|Y_s|^2$ on $\{|Y_s|>0\}$ and $0$ otherwise.
\end{proof}

\begin{corollary}[Deterministic moment bound]\label{cor:moment-bound}
For every $q\ge 1$ and every $t\ge0$,
\begin{equation}\label{eq:moment-direction}
\mathbb{E}\,|J_t(u)\xi|^{2q}
\le \exp\Big(q(2L+dL_\sigma^2)t + 2q^2 dL_\sigma^2 t\Big).
\end{equation}
Consequently,
\begin{equation}\label{eq:moment-op}
\mathbb{E}\,\|J_t(u)\|_{\mathrm{op}}^{2q}
\le \mathbb{E}\,\|J_t(u)\|_{\mathrm{HS}}^{2q}
\le d^{\,q}\exp\Big(q(2L+dL_\sigma^2)t + 2q^2 dL_\sigma^2 t\Big).
\end{equation}
\end{corollary}


\begin{definition}[Random Jacobian matrix/Lipschitz constant on a set]\label{def:jacobian-constant}
For a (nonempty) set $K\subset\mathbb{R}^d$, define
\begin{equation*}
\mathcal{K}_t(\omega;K):=\sup_{u\in K}\|J_t(u,\omega)\|_{\mathrm{op}}.
\end{equation*}
If $K$ is compact and $u\mapsto J_t(u,\omega)$ is continuous, then $\mathcal{K}_t(\omega;K)<\infty$.
\end{definition}

\begin{remark}[Why $\mathcal{K}_t(\omega;K)$ is finite and relation to global blow-up examples]
The counterexamples of Mohammed--Scheutzow show that even for globally Lipschitz coefficients one may fail to have a finite
global constant $\sup_{u\in\mathbb{R}^d}\|J_t(u,\omega)\|_{\mathrm{op}}$ for a fixed $(t,\omega)$.
In contrast, the quantity $\mathcal{K}_t(\omega;K)$ is taken over a \emph{compact} set $K$; under Assumption~\ref{ass:general_lipschitz}(iii),
for each fixed $(t,\omega)$ the map $u\mapsto J_t(u,\omega)$ is continuous, hence the supremum over $K$ is finite.
\end{remark}

\begin{corollary}[Pathwise LSI with the Jacobian matrix constant]\label{cor:lsi-jacobian-constant}
If $\mu_0$ satisfies $\mathrm{LSI}(C_0)$ and $\mathrm{supp}(\mu_0)\subset K$ for some compact $K$, then for every $t\ge0$ and a.s.\ $\omega$,
\begin{equation*}
\mu_t(\omega)=T_t(\omega)_{\#}\mu_0
\quad\Longrightarrow\quad
\mu_t(\omega)\text{ satisfies }\mathrm{LSI}\big(C_0\,\mathcal{K}_t(\omega;K)^2\big).
\end{equation*}
\end{corollary}

\bibliographystyle{amsplain}
\bibliography{bibliography}

\end{document}
